% Original analytic teaching diagram, 2026-11-10.
\definecolor{VB}{HTML}{7188A3}\definecolor{VBF}{HTML}{E5EAF0}
\definecolor{VG}{HTML}{829C8B}\definecolor{VGF}{HTML}{E8EEE9}
\definecolor{VY}{HTML}{B4A16F}\definecolor{VYF}{HTML}{F1ECDD}
\definecolor{VR}{HTML}{A7787D}\definecolor{VRF}{HTML}{F0E5E6}
\definecolor{VP}{HTML}{9588A5}\definecolor{VPF}{HTML}{ECE8F0}
\definecolor{VI}{HTML}{3B4252}\definecolor{VM}{HTML}{6B7280}\definecolor{VL}{HTML}{4C4D4F}
\begin{tikzpicture}[x=1mm,y=1mm,>=Stealth,font=\figurefont\fontsize{8}{10}\selectfont,
lab/.style={text=VM,font=\figurefont\fontsize{7.5}{9}\selectfont}]
% (a) gluing two local solutions
\draw[VL,line width=.75pt,->] (0,6)--(50,6) node[below,text=VI]{$t$};
\draw[VY,line width=2.4pt] plot[domain=20:45,samples=50,smooth] (\x,{16+7*sin(\x*9)});
\draw[VB,line width=1.1pt] plot[domain=5:32,samples=50,smooth] (\x,{16+7*sin(\x*9)});
\fill[VR] (12,{16+7*sin(108)}) circle (.85);
\node[text=VI] at (12,{16+7*sin(108)+3.4}) {$(t_0,\bm{x}_0)$};
\draw[VB,line width=1.1pt] (5,3.8)--(32,3.8); \node[lab,left] at (5,3.8) {$J_1$};
\draw[VY,line width=1.3pt] (20,1.8)--(45,1.8); \node[lab,right] at (45,1.8) {$J_2$};
\foreach \a in {20,32} \draw[VM,dash pattern=on .8pt off 1pt] (\a,1.5)--(\a,27);
\node[lab] at (26,27.8) {重叠};
\node[text=VI] at (25,-3.2) {(a)};
% (b) no first split
\begin{scope}[xshift=59mm]
\draw[VL,line width=.75pt,->] (0,6)--(50,6) node[below,text=VI]{$t$};
\fill[VBF,draw=VB,line width=.75pt] (14,10) rectangle (34,24);
\draw[VB,line width=1.3pt] plot[domain=3:44,samples=50,smooth] (\x,{13+.28*\x});
\draw[VR,line width=1.1pt,dash pattern=on 3pt off 1.8pt] plot[domain=24:40,samples=30,smooth] (\x,{13+.28*\x+.035*(\x-24)^2});
\fill[VI] (24,{13+.28*24}) circle (.75);
\draw[VM,dash pattern=on .8pt off 1pt] (24,6)--(24,{13+.28*24});
\node[text=VI,below] at (24,5.6) {$t^\ast$};
\node[lab] at (24,26.3) {局部唯一的窗};
\node[text=VR,right] at (36.5,27) {不能分叉};
\node[text=VI] at (25,-3.2) {(b)};
\end{scope}
\end{tikzpicture}
